\section{Learning}\label{sec:learning}

Now that the structure of a \gls{nn} has been detailed, the next step is to examine how its parameters are adjusted to accomplish the assigned task.
This brings the discussion to the core of \gls{dl}, namely the learning process itself.
As a reminder, the objective of training a \gls{nn} is to adjust its parameters so that the network provides a good approximation of the target function.
Mathematically, this objective is expressed as:
\begin{equation}
\min_{\theta} \frac{1}{N} \sum_{i=1}^{N} \mathcal{L}(\hat{f}_{\theta}(x^{(i)}), y^{(i)}).
\label{eq:training_objective}
\end{equation}

However, the \gls{nn} $\hat{f}$ used to approximate the target function $f$ is highly complex.
It involves a large number of parameters as well as many nonlinear components.
As a result, the optimum cannot be determined analytically by directly studying the properties of the function.
Instead, approximate optimization methods are required to find suitable parameter values for \gls{nn}.
The class of algorithms used for this optimization is known as gradient descent methods.
Numerous variants of gradient descent have been developed for \gls{nn} training, each designed to improve specific aspects of the optimization process.
For simplicity, the focus here is on \gls{sgd} with a batch size of 1.
As in previous cases, this choice is made because this setting captures the essential intuition behind the training method.

This section is organized into three steps.
The first derives the gradient of the loss using the chain rule.
The second shows how this gradient is used to update the network parameters.
The third discusses the convergence properties of the resulting algorithm.

\subsection{Computing the Gradient}

The first step of \gls{sgd} is to randomly initialize the parameter vector $\theta \in \mathbb{R}^P$ for a given network architecture $\hat{f}$.
This initial state is denoted by $\theta^{(1)}$.
The goal is to iteratively update this vector so that the network gets better at approximating the target function.
To do this, the influence of each parameter $w_k$ on the loss $\mathcal{L}(\hat{f}_{\theta}(x^{(i)}), y^{(i)})$ must be quantified, which is given by the partial derivative $\frac{\partial \mathcal{L}(\hat{f}_{\theta}(x^{(i)}), y^{(i)})}{\partial w_k}$.

However, $w_k$ can belong to any layer of the network.
Computing this derivative therefore requires propagating the effect of a change in $w_k$ through all the layers that follow it, up to the loss.
To do so, each layer is considered as a function of several variables, whose inputs are collected in the vector $\mathbf{h}^{(l-1)}$ and whose outputs are collected in the vector $\mathbf{h}^{(l)}$.
Since these quantities depend on several variables, the Jacobian matrix is used as a generalization of the derivative.
The Jacobian matrix quantifies how a small change in each input component affects each output component.
For a function $\mathbf{f}\colon\mathbb{R}^n \rightarrow \mathbb{R}^m$, $J_{\mathbf{f}}(\mathbf{x})$ denotes the Jacobian of $\mathbf{f}$ evaluated at $\mathbf{x}$, which collects the partial derivatives of every output component with respect to every input component:
\begin{equation}
J_{\mathbf{f}}(\mathbf{x}) =
\begin{bmatrix}
\frac{\partial f_1}{\partial x_1} & \cdots & \frac{\partial f_1}{\partial x_n} \\
\vdots & \ddots & \vdots \\
\frac{\partial f_m}{\partial x_1} & \cdots & \frac{\partial f_m}{\partial x_n}
\end{bmatrix}.
\label{eq:jacobian}
\end{equation}
Each entry indicates how one output component reacts to a small change in one input component.

To propagate the effect of $w_k$ through the successive layers, the derivative of a composition of several functions is required.
For a composition of differentiable functions $\mathbf{f} \circ \mathbf{g}$, the chain rule states that its Jacobian, evaluated at $\mathbf{x}$, is the product of the Jacobians of each function, each evaluated at the appropriate point:
\begin{equation}
J_{\mathbf{f} \circ \mathbf{g}}(\mathbf{x}) = J_{\mathbf{f}}(\mathbf{g}(\mathbf{x})) \, J_{\mathbf{g}}(\mathbf{x}).
\label{eq:chain_rule_generic}
\end{equation}

This is the tool required to determine the impact of $w_k$ on the output of the network, regardless of where it is located.
A network computes its output through a sequence of layer activations $\mathbf{h}^{(1)}, \dots, \mathbf{h}^{(L)}$, with each layer depending only on the one before it.
As a result, a parameter $w_k$ belonging to layer $l$ affects the loss $\mathcal{L}(\hat{f}_{\theta}(x^{(i)}), y^{(i)})$ through a chain of dependencies, it first changes $\mathbf{h}^{(l)}$, which changes $\mathbf{h}^{(l+1)}$, and so on until the output layer $\mathbf{h}^{(L)}$, from which the loss is computed.
Since $x^{(i)}$ and $y^{(i)}$ are fixed for a given example, the loss can be viewed as a function of $w_k$ alone, and the effect of a small change in $w_k$ on this value is sought.
Applying the chain rule across this chain gives
\begin{equation}
\frac{\partial \mathcal{L}(\hat{f}_{\theta}(x^{(i)}), y^{(i)})}{\partial w_k}
=
J_{\mathcal{L}}(\mathbf{h}^{(L)}, y^{(i)})
\, J_{\mathbf{h}^{(L)}}(\mathbf{h}^{(L-1)})
\cdots
J_{\mathbf{h}^{(l+1)}}(\mathbf{h}^{(l)})
\, J_{\mathbf{h}^{(l)}}(w_k),
\label{eq:chain_rule}
\end{equation}
where:
\begin{itemize}
\item $J_{\mathcal{L}}(\mathbf{h}^{(L)}, y^{(i)})$ denotes the Jacobian of $\mathcal{L}$ with respect to its first argument, evaluated at $(\mathbf{h}^{(L)}, y^{(i)})$,
\item each subsequent $J_{\mathbf{h}^{(j+1)}}(\mathbf{h}^{(j)})$ denotes the Jacobian of layer $j+1$ with respect to its input $\mathbf{h}^{(j)}$,
\item $J_{\mathbf{h}^{(l)}}(w_k)$ denotes the Jacobian of layer $l$ with respect to the parameter $w_k$.
\end{itemize}
Each factor is therefore the Jacobian of one layer with respect to its input.
Multiplying these local Jacobians together propagates the effect of $w_k$ all the way to the loss.
This layer-by-layer propagation is referred to as backpropagation.

Once $\frac{\partial \mathcal{L}(\hat{f}_{\theta}(x^{(i)}), y^{(i)})}{\partial w_k}$ can be computed for every parameter, all $P$ of these values are stacked into one vector, called the gradient of the loss with respect to the parameters:
\begin{equation}
\nabla_{\theta} \mathcal{L}(\hat{f}_{\theta}(x^{(i)}), y^{(i)}) =
\left[
\frac{\partial \mathcal{L}(\hat{f}_{\theta}(x^{(i)}), y^{(i)})}{\partial w_1},
\dots,
\frac{\partial \mathcal{L}(\hat{f}_{\theta}(x^{(i)}), y^{(i)})}{\partial w_P}
\right].
\label{eq:gradient}
\end{equation}
Each component answers the same question for a different parameter, namely, if this parameter is slightly increased, does the loss go up or down, and by how much.
A positive value means increasing the parameter increases the loss, whereas a negative value means the opposite.
The gradient thus indicates, for every parameter, which direction to move in order to reduce the loss.

\subsection{Update rule}

Now that information is available on how changes in the parameters affect the loss, the objective is to minimize it.
To do so, the parameters are updated in the direction opposite to the gradient.
Indeed, the gradient is followed in the opposite direction because it indicates the direction in which the loss increases, whereas the objective is to reduce it.
This leads to the gradient descent update rule:
\begin{equation}
\theta^{(n+1)} = \theta^{(n)} - \eta \nabla_{\theta^{(n)}} \mathcal{L}(\hat{f}_{\theta^{(n)}}(x^{(i)}), y^{(i)}),
\label{eq:gradient_descent_update}
\end{equation}
where $\eta > 0$ is the learning rate, which controls the step size.

This parameter $\eta$ must remain small, as the gradient only provides reliable information in a local neighborhood of the current parameters.
As a result, a single step is not sufficient to significantly reduce the loss, so the process is iterated many times.
Producing a sequence of parameter vectors $\theta_1, \theta_2, \dots$ that progressively improve the model.
In practice, training is stopped when successive updates no longer produce a significant decrease in the loss.
Algorithm~\ref{alg:stochastic_gradient_descent} summarizes this whole training procedure.

\begin{algorithm}[t]
\caption{Stochastic Gradient Descent}
\label{alg:stochastic_gradient_descent}

\Input{
\\
Dataset $\mathcal{D} = \{(x^{(i)}, y^{(i)})\}_{i=1}^{N}$\\
Architecture $\hat{f}$\\
Loss function $\mathcal{L}$\\
Learning rate $\eta > 0$\\
Initial parameters $\theta$ (optional)
}

\Output{
\\
Trained parameters $\theta$
}

\AlgoRule

\uIf{$\theta$ is given}{
    $\theta^{(1)} \leftarrow \theta$\;
}\uElse{
    Initialize $\theta^{(1)}$ randomly\;
}

\BlankLine
$n \leftarrow 1$\;
\While{stopping criterion not satisfied}{
    Select a random example $(x^{(i)}, y^{(i)}) \in \mathcal{D}$\;

    \BlankLine
    \tcp{Model prediction for the selected example}
    $\hat{y}^{(i)} \leftarrow \hat{f}_{\theta^{(n)}}(x^{(i)})$\;

    \BlankLine
    \tcp{Per-example loss between the prediction and the ground truth}
    $\loss{\theta^{(n)}}{i} \leftarrow \mathcal{L}(\hat{y}^{(i)}, y^{(i)})$\;

    \BlankLine
    \tcp{Gradient of the per-example loss, since it is differentiable}
    $g \leftarrow \nabla_{\theta} \loss{\theta^{(n)}}{i}$\;

    \BlankLine
    \tcp{Error minimization via a gradient descent step}
    $\theta^{(n+1)} \leftarrow \theta^{(n)} - \eta \, g$\;

    $n \leftarrow n + 1$\;
}

\BlankLine

\Return{$\theta^{(n)}$}

\end{algorithm}


To develop an intuition for this update rule, Figure~\ref{fig:simple_gradient_descent} illustrates it applied to the simple convex loss function $\mathcal{L}(w_1,w_2)=w_1^2+w_2^2$, where the red path shows the parameter updates converging to the minimum.
The point $\theta_1$ represents the initial parameter values, with $w_1$ and $w_2$ randomly initialized, and the subsequent points, from $\theta_2$ to $\theta_6$, represent the successive updates produced by applying Equation~\ref{eq:gradient_descent_update}, progressively moving the trajectory toward the minimum of the loss function.

% % ---------- Customisation ----------
% Arrow style
\def\arrowcolor{red}
\def\arrowtype{->}          % -stealth, -latex, etc.
\def\arrowthickness{thick}

% Label style
\def\labelfont{\footnotesize}
\def\labelfill{white}
\def\labelprefix{\theta}

% Label offset (in axis units)
\def\labeloffsetx{0}
\def\labeloffsety{0.75}
\def\labeloffsetz{-10}
% -----------------------------------

\begin{figure}[t]
    \centering
    \resizebox{0.85\textwidth}{!}{
        \begin{tikzpicture}
        \begin{axis}[
                % title={$\mathcal{L}(\theta) = w_{1}^2 + w_{2}^2$},
                view={-60}{60},
                xlabel={$w_{1}$},
                ylabel={$w_{2}$},
                zlabel={$\mathcal{L}(w_{1}, w_{2})$},
                colormap/viridis,
                % colorbar,
                % colorbar style={
                %     y dir=reverse,
                %     yticklabel={\pgfmathparse{-\tick}\pgfmathprintnumber{\pgfmathresult}},
                % },
                mesh/ordering=y varies,
                scaled z ticks=false,
        ]

        % ---------- Loss surface ----------
        \addplot3[
            surf,
            shader=flat,
            opacity=0.8,
            faceted color=none,
            point meta=-z,
        ] table {figures/deepl_learning/gradient_descent/simple_grandient_descent/loss_surface.data};

        % ---------- Optimization path ----------
        \pgfplotstableread{figures/deepl_learning/gradient_descent/simple_grandient_descent/optimization_path.data}\datatable
        \pgfplotstablegetrowsof{\datatable}
        \pgfmathtruncatemacro{\numpts}{\pgfplotsretval}
        \pgfmathtruncatemacro{\lastpoint}{\numpts-1}
        \pgfmathtruncatemacro{\lastseg}{\numpts-2}

        % 1. Draw every segment with an arrow
        \foreach \i in {0,...,\lastseg}{
            \pgfmathtruncatemacro{\j}{\i+1}
            \pgfplotstablegetelem{\i}{0}\of\datatable \let\xstart\pgfplotsretval
            \pgfplotstablegetelem{\i}{1}\of\datatable \let\ystart\pgfplotsretval
            \pgfplotstablegetelem{\i}{2}\of\datatable \let\zstart\pgfplotsretval
            \pgfplotstablegetelem{\j}{0}\of\datatable \let\xend\pgfplotsretval
            \pgfplotstablegetelem{\j}{1}\of\datatable \let\yend\pgfplotsretval
            \pgfplotstablegetelem{\j}{2}\of\datatable \let\zend\pgfplotsretval
            \addplot3[
                \arrowcolor,
                \arrowtype,
                \arrowthickness,
                mark=none,
            ] coordinates {
                (\xstart,\ystart,\zstart) (\xend,\yend,\zend)
            };
        }

        % 2. Place a label on every point
        \foreach \k in {0,...,\lastpoint}{
            \pgfmathtruncatemacro{\labelidx}{\k+1}
            \pgfplotstablegetelem{\k}{0}\of\datatable \let\x\pgfplotsretval
            \pgfplotstablegetelem{\k}{1}\of\datatable \let\y\pgfplotsretval
            \pgfplotstablegetelem{\k}{2}\of\datatable \let\z\pgfplotsretval
            % Expand coordinates immediately, but protect font commands
            \edef\templabel{
                \noexpand\node[
                    font=\noexpand\labelfont,
                    fill=\noexpand\labelfill,
                    inner sep=1pt,
                ] at (axis cs: \x+\labeloffsetx,\y+\labeloffsety,\z+\labeloffsetz)
                    {$\labelprefix_{\labelidx}$};
            }
            \templabel
        }

        \end{axis}
        \end{tikzpicture}
    }
    \caption{Gradient descent in a case $\mathcal{L}(w_{1}, w_{2}) = w_{1}^2 + w_{2}^2$.}
    \label{fig:simple_gradient_descent}
\end{figure}

\subsection{Convergence Properties}

Having established how \gls{sgd} updates the parameters, a natural question is whether this procedure is guaranteed to converge.
It has been mathematically shown that gradient descent applied to a convex loss with respect to the parameters converges to a global optimum, independently of the initialization.
However, \gls{nn} training involves highly non-convex loss surfaces, so the optimization trajectory becomes strongly dependent on the initial conditions, and no such guarantee holds.
Figure~\ref{fig:multiple_gradient_descents_non_convex} illustrates the effect of non-convex functions, different starting points lead to distinct local minima.

% % --- General path for the data files ---
\def\datapath{figures/deepl_learning/gradient_descent/multiple_gradient_descents_non_convex/}
% ----------------------------------------

\begin{figure}[t]
    \centering
    % First row: Goldstein-Price and Rosenbrock
    \begin{subfigure}[b]{0.48\textwidth}
        \centering
        \resizebox{\textwidth}{!}{%
        \begin{tikzpicture}
        \begin{axis}[
            view={-45}{45},
            colormap/viridis,
            mesh/ordering=y varies,
            scaled z ticks=false,
            xticklabels={}, yticklabels={}, zticklabels={},
        ]
            \addplot3[
                surf, shader=flat, opacity=0.8,
                faceted color=none, point meta=-z,
            ] table {\datapath loss_surface_goldstein.data};
            \addplot3[red, line width=1.2pt] table {\datapath optimization_path_goldstein_a.data};
            \addplot3[red, line width=1.2pt] table {\datapath optimization_path_goldstein_b.data};
        \end{axis}
        \end{tikzpicture}%
        }
        \caption{Goldstein-Price}
        \label{fig:goldstein}
    \end{subfigure}
    \hfill
    \begin{subfigure}[b]{0.48\textwidth}
        \centering
        \resizebox{\textwidth}{!}{%
        \begin{tikzpicture}
        \begin{axis}[
            view={-45}{-45},
            colormap/viridis,
            mesh/ordering=y varies,
            scaled z ticks=false,
            xticklabels={}, yticklabels={}, zticklabels={},
        ]
            \addplot3[
                surf, shader=flat, opacity=0.8,
                faceted color=none, point meta=-z,
            ] table {\datapath loss_surface_rosenbrock.data};
            \addplot3[red, line width=1.2pt] table {\datapath optimization_path_rosenbrock_a.data};
            \addplot3[red, line width=1.2pt] table {\datapath optimization_path_rosenbrock_b.data};
        \end{axis}
        \end{tikzpicture}%
        }
        \caption{Rosenbrock}
        \label{fig:rosenbrock}
    \end{subfigure}

    \vspace{1em}

    % Second row: Sombrero and Ackley
    \begin{subfigure}[b]{0.48\textwidth}
        \centering
        \resizebox{\textwidth}{!}{%
        \begin{tikzpicture}
        \begin{axis}[
            view={-45}{45},
            colormap/viridis,
            mesh/ordering=y varies,
            scaled z ticks=false,
            xticklabels={}, yticklabels={}, zticklabels={},
        ]
            \addplot3[
                surf, shader=flat, opacity=0.8,
                faceted color=none, point meta=-z,
            ] table {\datapath loss_surface_sombrero.data};
            \addplot3[red, line width=1.2pt] table {\datapath optimization_path_sombrero_a.data};
            \addplot3[red, line width=1.2pt] table {\datapath optimization_path_sombrero_b.data};
        \end{axis}
        \end{tikzpicture}%
        }
        \caption{Sombrero}
        \label{fig:sombrero}
    \end{subfigure}
    \hfill
    \begin{subfigure}[b]{0.48\textwidth}
        \centering
        \resizebox{\textwidth}{!}{%
        \begin{tikzpicture}
        \begin{axis}[
            view={-135}{45},
            colormap/viridis,
            mesh/ordering=y varies,
            scaled z ticks=false,
            xticklabels={}, yticklabels={}, zticklabels={},
        ]
            \addplot3[
                surf, shader=flat, opacity=0.8,
                faceted color=none, point meta=-z,
            ] table {\datapath loss_surface_ackley.data};
            \addplot3[red, line width=1.2pt] table {\datapath optimization_path_ackley_a.data};
            \addplot3[red, line width=1.2pt] table {\datapath optimization_path_ackley_b.data};
        \end{axis}
        \end{tikzpicture}%
        }
        \caption{Ackley}
        \label{fig:ackley}
    \end{subfigure}

    \caption{Comparison of loss surfaces and gradient descent trajectories.}
    \label{fig:multiple_gradient_descents_non_convex}
\end{figure}


Despite the lack of theoretical guarantees of convergence to a global optimum, the simple procedure of random initialization followed by iterative gradient-based updates is sufficient in practice to train complex \gls{nn}.